\chapter{Design details}
\label{app:design}

This appendix gives the definitions and properties that the chapters rely on.

\section{Tasks and environments}

\paragraph{Rows.}
Features are drawn as $x \sim \mathcal{N}(0, I_{10})$, the environment's perturbation is applied, the clean label is computed from the load-bearing features $C$, the shortcut $x_8$ is generated from the clean label, $x_9 \sim \mathcal{N}(0,1)$, and labels are flipped with probability $\eta = 0.02$.
Features are standardised with the pool's statistics and shown in a random order $\pi_t$ per task, so no role can be read from a position.

\paragraph{Rules.}
A \emph{linear} task has $y_{\mathrm{clean}} \sim \mathrm{Bernoulli}\big(\sigma(c^\top x_C)\big)$ with $c_j = u_j r_j$, $u_j \sim \mathrm{U}(1,2)$, $r_j \in \{\pm 1\}$; $P(y_{\mathrm{clean}} = 1 \mid x)$ is strictly monotone in each $x_j$, so feature $j$ has the direction $s_j = \operatorname{sign} c_j$.
A \emph{tree} task splits each of three features at zero and labels the eight leaves; the leaf tables allowed (complement-antisymmetric, neither a dictator nor parity) are exactly the signed majorities $y_{\mathrm{clean}} = 1\big[\mathrm{Maj}(s_1 x_1, s_2 x_2, s_3 x_3) > 0\big]$, so tree tasks are monotone too.

\paragraph{The shortcut.}
With $d = \Phi^{-1}(s)$, $x_8 = \sigma_t\big((2y_{\mathrm{clean}} - 1)d + \varepsilon\big)$, $\varepsilon \sim \mathcal{N}(0,1)$.

\begin{proposition}[Sign agreement]
$P\big(\operatorname{sign}(\sigma_t x_8) = \operatorname{sign}(2y_{\mathrm{clean}}-1)\big) = s$; under spurious reversal ($d \to -d$) it is $1 - s$.
\end{proposition}
\begin{proof}
For $y = 1$, $P(\sigma_t x_8 > 0) = P(Z > -d) = \Phi(d) = s$; $y = 0$ is symmetric. Reversal replaces $d$ by $\Phi^{-1}(1-s) = -d$.
\end{proof}

\paragraph{Co-variate shift.}
Two load-bearing features $i, j$ are moved in opposite label directions and a distractor $d \notin C$ at random:
$\delta_i = s_i m/|c_i|$, $\delta_j = -s_j m/|c_j|$, $\delta_d \sim \pm\mathrm{U}(1,2)$, with the push $m$ chosen so that every shifted feature moves by 1--2 standard deviations (for trees $|c| = 1$, $m \sim \mathrm{U}(1,2)$).

\begin{lemma}[The label rate is unchanged]
(a) For linear tasks, $c^\top(x + \delta) = c^\top x$ for every row. (b) For tree tasks, $P(y_{\mathrm{clean}} = 1)$ stays $1/2$.
\end{lemma}
\begin{proof}
(a) $c_i\delta_i = m$, $c_j\delta_j = -m$, and the distractor carries no weight. (b) The three votes are independent. After the shift the two moved votes are 1 with probabilities $p_i = \Phi(m)$ and $p_j = 1 - \Phi(m)$, and the third with $\tfrac12$, so the majority is 1 with probability $p_i p_j + \tfrac12[p_i(1-p_j) + (1-p_i)p_j] = \tfrac12(p_i + p_j) = \tfrac12$.
\end{proof}

\section{Naming}

Each domain has antonym pairs $(p, q)$, where $p$ is associated in everyday language with label 1 and $q$ with label 0, and a set of weak-prior names.
Each load-bearing feature $j$ is assigned a pair: \emph{aligned} names it $p$ if $s_j = +1$ and $q$ if $s_j = -1$; \emph{flipped} uses the other member.
Every other feature gets a weak-prior name, the same in both conditions.

\begin{proposition}[Identification]
For a fixed task, demonstration set and query, the aligned and flipped prompts differ only in the names of the three load-bearing features. A predictor that ignores names has the same expected accuracy under both, so a difference in the model's accuracy is attributable to the meaning of the names.
\end{proposition}

\section{Counter-spurious selection}

Balancing the four (label, shortcut agrees) cells is not enough: agreeing rows lie further from zero ($\mathbb{E}[|x_8| \mid \text{agree}] = d + \phi(d)/\Phi(d) > \mathbb{E}[|x_8| \mid \text{disagree}] = \phi(d)/(1-\Phi(d)) - d$), so $x_8$ still correlates with the label (about 0.33 at $s = 0.85$).
\strat{counter\_spurious} therefore takes, within each label, $k/4$ clean disagreeing rows and pairs each with the clean agreeing row of the same label nearest in $|x_8|$ \citep{rubin1973matching,stuart2010matching}.
With perfect matching, sign agreement is $\tfrac12$ and $\mathrm{Cov}(x_8, 2y-1) = 0$; on the evaluation tasks the implementation gives $\mathrm{corr}(x_8, y) = 0.005$.
\sata's matched target (\autoref{sec:sata-targets}) is each row's probability $\pi_i$ of appearing in such a set built for the query, with rows in the query's regime (the sign pattern of its three load-bearing features) taken first, normalized: $q_i = \pi_i / \sum_j \pi_j$; label-noise rows have $\pi_i = 0$, and nothing reads the query's label.
\sata is trained with the loss $\mathrm{KL}(q \,\Vert\, p_\theta)$.

\section{Scoring and caching}

\begin{lemma}[Margin]
With $p_1 = e^{\ell_1}/(e^{\ell_0} + e^{\ell_1})$ and $m = \ell_1 - \ell_0$, $p_1 = \sigma(m)$, and $m$ is the same whether computed from logits or from full-vocabulary log-probabilities.
\end{lemma}
\begin{proof}
$p_1/(1 - p_1) = e^{\ell_1 - \ell_0}$; log-probabilities subtract a common log-sum-exp, which cancels.
\end{proof}

\begin{proposition}[Prefix caching is exact]
For a decoder-only transformer with causal attention, running the prefix $P$ once and each suffix $S$ against its cached keys and values gives the same final-position logits as a full pass on $P \Vert S$, provided $\mathrm{tok}(P \Vert S) = \mathrm{tok}(P) \Vert \mathrm{tok}(S)$.
\end{proposition}
\begin{proof}
States at positions $< |P|$ depend only on those positions' tokens (attention reads positions $\le t$), so appending $S$ does not change them, and the states of $S$ are computed from the same inputs either way. The runner checks the tokenisation condition for every suffix.
\end{proof}

\begin{lemma}[Logit lens at the last layer]
With $w = W_U[\texttt{1}] - W_U[\texttt{0}]$ and $h_L$ the final (normalized) state, $m_L = w^\top h_L = \ell_1 - \ell_0 = m$.
\end{lemma}

\section{Statistics}

\paragraph{Cells and contrasts.}
\AUROC is computed within each cell (task, strategy, seed, naming, environment) over its 20 label-balanced queries and averaged over seeds.
A contrast is a per-task difference of such values (for example, strategy $A$ minus strategy $B$ on the same queries, or ID minus a shifted environment), averaged over the 24 tasks.

\paragraph{Hierarchical bootstrap.}
Each of $B = 2000$ draws resamples tasks with replacement and, within each resampled task, its positive and negative queries separately; the percentile interval of the resampled estimates is reported, with one-sided $p$-values from the share of draws on the wrong side of zero.
Because tasks differ in their rules, the variance of the mean is dominated by the between-task term ($\sigma_T^2/T + \sigma_W^2/(Tn)$; \citealp{cochran1977sampling}), so tasks are the unit of resampling.

\paragraph{Multiplicity.}
Holm's step-down correction \citep{holm1979simple} is applied within each pre-registered family: ten contrasts for \RQ1--\RQ3 and six for \sata.
\autoref{tab:contrasts} lists every confirmatory estimate.

\begin{table}[h]
\centering\small
\caption{Every confirmatory contrast (Qwen2.5-7B-Instruct, 24 tasks; one-sided hierarchical bootstrap, Holm within each family). All are differences in \AUROC except demo-following (in the demo-following index) and feature reliance (in Spearman $\rho$). Bold: significant after correction.}
\label{tab:contrasts}
\begin{tabular}{@{}lrrr@{}}
\toprule
\textbf{Contrast} & \textbf{Estimate [95\% CI]} & \textbf{Holm $p$} & \textbf{Tasks} \\
\midrule
\multicolumn{4}{@{}l}{\emph{Shift, names and the prior (\RQ1--\RQ3)}} \\
Covariate shift: ID $-$ covariate (random, abstract) & $+0.024$ [$-0.044$, 0.093] & 0.95 & 15/24 \\
Shortcut reversal: ID $-$ reversal (random, abstract) & $+0.115$ [0.017, 0.218] & 0.063 & 16/24 \\
\textbf{Flipped names: aligned $-$ flipped (random, ID)} & $\mathbf{+0.204}$ [0.131, 0.280] & \textbf{0.005} & 22/24 \\
counter\_spurious $-$ label\_diversity, reversal & $+0.035$ [$-0.019$, 0.086] & 0.59 & 16/24 \\
similarity $-$ label\_diversity, covariate & $-0.029$ [$-0.110$, 0.048] & 1.00 & 11/24 \\
feature\_range $-$ label\_diversity, covariate & $-0.078$ [$-0.141$, $-0.018$] & 1.00 & 6/24 \\
counter\_prior $-$ label\_diversity, flipped, ID & $+0.023$ [$-0.015$, 0.064] & 0.70 & 7/24 \\
\textbf{counter\_spurious's gain: reversal $-$ covariate} & $\mathbf{+0.111}$ [0.033, 0.196] & \textbf{0.024} & 17/24 \\
\textbf{Demo-following index: abstract $-$ flipped} & $\mathbf{+0.123}$ [0.079, 0.175] & \textbf{0.005} & 23/24 \\
Feature reliance $\rho(\pi_{\mathrm{true}}, \pi_{\mathrm{behav}})$: abstract $-$ flipped & $+0.026$ [$-0.141$, 0.174] & 1.00 & 13/24 \\
\midrule
\multicolumn{4}{@{}l}{\emph{\sata (\RQ4)}} \\
\sata{} $-$ label\_diversity, reversal & $+0.032$ [$-0.024$, 0.089] & 0.79 & 14/24 \\
\sata{} $-$ label\_diversity, covariate & $-0.096$ [$-0.154$, $-0.041$] & 1.00 & 5/24 \\
\sata{} $-$ label\_diversity, ID & $-0.040$ [$-0.102$, 0.023] & 1.00 & 8/24 \\
\sata{} $-$ similarity, mean of 3 environments & $-0.020$ [$-0.064$, 0.026] & 1.00 & 11/24 \\
\sata{} $-$ \sata{} without the query, mean of 3 & $-0.025$ [$-0.063$, 0.010] & 1.00 & 10/24 \\
\sata{} $-$ label\_diversity, flipped, ID & $-0.014$ [$-0.077$, 0.048] & 1.00 & 13/24 \\
\bottomrule
\end{tabular}
\end{table}

\section{Reproducibility}
\label{app:repro}

The code and configuration and the pre-registered contrast families all are available at \url{https://github.com/Chenuka-Garusinghe/LLM-ICL-OOD-Honours/tree/synthetic-experiments/sata-project}, with tests covering the invariants of the generator, the selection engine, the model runner and the probes.

Every prediction is stored with its demonstration rows, order and prompt, so every number and figure here can be recomputed from the result files.